% ==============================================================================
% Chapitre 3 — Composants Applicatifs
% Dossier d'Architecture LockBits
% ==============================================================================

\chapter{Composants Applicatifs}

Ce chapitre décrit en détail les quatre composants applicatifs du système
LockBits : le serveur EDR, l'agent EDR, le site client et le chatbot.
Pour chaque composant, nous présentons l'architecture interne, les
technologies utilisées, les fonctionnalités principales et les modalités
de déploiement.


% ==============================================================================
\section{EDR Server}
% ==============================================================================

Le serveur EDR (Endpoint Détection and Response) constitue le c\oe{}ur
central de la plateforme. Il reçoit, stocke et analyse les événements de
sécurité remontés par les agents déployés sur les postes clients.

% ------------------------------------------------------------------------------
\subsection{Architecture du serveur}
% ------------------------------------------------------------------------------

Le serveur est développé en \textbf{Python 3.12} avec le framework
\textbf{FastAPI}. Il repose sur une architecture modulaire ou chaque
fichier du répertoire \texttt{server/} remplit une responsabilité unique :

\begin{itemize}
  \item \texttt{main.py} -- point d'entrée de l'application FastAPI,
    configuration CORS, instrumentation Prometheus ;
  \item \texttt{api.py} -- routeur contenant tous les endpoints REST ;
  \item \texttt{database.py} -- configuration SQLAlchemy et gestion des
    sessions ;
  \item \texttt{models.py} -- modèles ORM (Agent, Event, HashCache) ;
  \item \texttt{schemas.py} -- schémas Pydantic pour la validation des
    données ;
  \item \texttt{config.py} -- chargement de la configuration depuis les
    variables d'environnement ;
  \item \texttt{vt\_worker.py} -- worker asynchrone d'enrichissement
    VirusTotal ;
  \item \texttt{glpi\_client.py} -- client de création de tickets GLPI ;
  \item \texttt{metrics.py} -- métriques Prometheus personnalisées ;
  \item \texttt{ratelimit.py} -- configuration du rate limiting avec
    slowapi ;
  \item \texttt{storage.py} -- couche d'accès aux données.
\end{itemize}

Le démarrage du serveur s'effectue avec \texttt{uvicorn}, un serveur ASGI
performant, via la commande suivante :

\begin{lstlisting}[language=bash, caption={Lancement du serveur EDR}, label={lst:edr-launch}]
python -m uvicorn server.main:app --host 0.0.0.0 --port 8000
\end{lstlisting}

Le point d'entrée \texttt{main.py} initialise la base de données au
démarrage via le gestionnaire de cycle de vie \texttt{lifespan} de
FastAPI. Il configure le middleware CORS, le rate limiting, et expose
automatiquement un endpoint \texttt{/metrics} pour Prometheus. La
documentation OpenAPI est accessible à \texttt{/docs} (Swagger UI) sans
aucune configuration supplémentaire.

% ------------------------------------------------------------------------------
\subsection{API REST}
% ------------------------------------------------------------------------------

Le routeur défini dans \texttt{api.py} expose les endpoints suivants :

\begin{itemize}
  \item \texttt{GET /health} -- healthcheck du serveur (verification base
    de données incluse) ;
  \item \texttt{GET /api/v1/stats} -- statistiques globales (nombre
    d'agents, d'événements, alertes VT) ;
  \item \texttt{GET /api/v1/agents} -- liste des agents enregistrés ;
  \item \texttt{GET /api/v1/agents/\{agent\_id\}} -- détail d'un agent ;
  \item \texttt{GET /api/v1/events} -- événements avec filtres optionnels
    (\texttt{agent\_id}, \texttt{event\_type}, \texttt{severity},
    \texttt{hostname}) et pagination (\texttt{page}, \texttt{page\_size}) ;
  \item \texttt{GET /api/v1/events/\{event\_id\}} -- détail d'un événement ;
  \item \texttt{POST /api/v1/events} -- ingestion d'un événement
    (authentifié par \texttt{AUTH\_TOKEN}) ;
  \item \texttt{POST /api/v1/events/batch} -- ingestion par lot ;
  \item \texttt{POST /api/v1/heartbeat} -- heartbeat des agents.
\end{itemize}

L'extrait suivant montre la simplicité d'un endpoint FastAPI :

\begin{lstlisting}[language=Python, caption={Endpoint de healthcheck}, label={lst:health-endpoint}]
@router.get("/health", response_model=HealthResponse)
def health_check(db: Session = Depends(get_db)) -> HealthResponse:
    uptime = time.time() - _startup_time

    try:
        db.execute(text("SELECT 1"))
        db_ok = True
        db_health.set(1)
    except Exception:
        db_ok = False
        db_health.set(0)

    return HealthResponse(
        status="ok",
        version="1.0.0",
        db_ok=db_ok,
        uptime=uptime,
        timestamp=datetime.utcnow(),
    )
\end{lstlisting}

Tous les endpoints d'écriture sont protégés par une authentification
Bearer token. Le token est vérifié par la fonction \texttt{verify\_auth}
qui compare le header \texttt{Authorization} avec la variable
d'environnement \texttt{AUTH\_TOKEN}. Lorsque \texttt{AUTH\_TOKEN} n'est
pas définie, l'authentification est désactivée (mode développement).

% ------------------------------------------------------------------------------
\subsection{Base de données}
% ------------------------------------------------------------------------------

Le serveur utilise \textbf{SQLAlchemy} comme ORM et supporte deux moteurs
de base de données :

\begin{itemize}
  \item \textbf{PostgreSQL} en production via la variable
    \texttt{DATABASE\_URL} ;
  \item \textbf{SQLite} possible en développement (fichier local \texttt{miniedr.db}).
\end{itemize}

Le schéma comprend trois tables principales :

\begin{itemize}
  \item \texttt{agents} -- stocke les métadonnées des agents (ID,
    hostname, plateforme, version, compagnie, date de premiere et
    dernière connexion, statut) ;
  \item \texttt{events} -- stocke les événements de sécurité (type,
    sévérité, payload JSON, tags JSON, état d'enrichissement VT) ;
  \item \texttt{hash\_cache} -- cache des résultats VirusTotal pour
    éviter les appels API répétés.
\end{itemize}

\begin{lstlisting}[language=Python, caption={Modele SQLAlchemy pour les agents}, label={lst:agent-model}]
class Agent(Base):
    """Registered agent metadata and status."""
    __tablename__ = "agents"

    id = Column(Integer, primary_key=True)
    agent_id = Column(String(36), unique=True, nullable=False, index=True)
    hostname = Column(String(255), nullable=False)
    platform = Column(String(50), nullable=False)
    agent_version = Column(String(50), default="")
    company = Column(String(255), default="", index=True)
    first_seen = Column(DateTime, default=datetime.utcnow, nullable=False)
    last_seen = Column(DateTime, default=datetime.utcnow,
                       onupdate=datetime.utcnow, nullable=False)
    status = Column(String(20), default="unknown")
    ip_address = Column(String(45), default="")
\end{lstlisting}

% ------------------------------------------------------------------------------
\subsection{Integration VirusTotal}
% ------------------------------------------------------------------------------

Le worker VirusTotal (\texttt{vt\_worker.py}) est un composant asynchrone
qui enrichit les événements avec des informations de réputation sur les
fichiers suspects. Lorsqu'un événement FIM contient un hash SHA256, le
worker interroge l'API publique de VirusTotal.

Les appels sont limités à 4 requêtes par minute via un \texttt{Semaphore}
asyncio, ce qui respecte les quotas de l'API gratuite. Un cache local
(\texttt{hash\_cache}) evite de interroger VirusTotal plusieurs fois pour
le meme hash.

Le worker est active via les variables d'environnement
\texttt{VT\_ENABLED} et \texttt{VT\_API\_KEY}.

% ------------------------------------------------------------------------------
\subsection{Integration GLPI}
% ------------------------------------------------------------------------------

L'intégration avec GLPI (\texttt{glpi\_client.py}) permet de créer
automatiquement un ticket GLPI pour chaque événement de sécurité accepté
par le serveur. Cette fonctionnalité est particulièrement utile pour les
équipes SOC qui utilisent GLPI comme outil de ticketing.

Le client utilise l'API REST v2 de GLPI avec authentification OAuth2
(grant type \texttt{password}). La configuration nécessite les variables
suivantes :

\begin{itemize}
  \item \texttt{GLPI\_ENABLED} -- active ou desactive l'intégration ;
  \item \texttt{GLPI\_API\_URL} -- URL de l'API GLPI ;
  \item \texttt{GLPI\_OAUTH\_CLIENT\_ID} et
    \texttt{GLPI\_OAUTH\_CLIENT\_SECRET} -- identifiants OAuth2 ;
  \item \texttt{GLPI\_API\_USERNAME} et \texttt{GLPI\_API\_PASSWORD} --
    credentials de l'API.
\end{itemize}

Le ticket créé résume l'événement avec sa sévérité, son type, le hostname
concerné, et le hash du fichier suspect. Le demandeur du ticket est
associé à l'utilisateur GLPI correspondant à l'email transmis par l'agent.

% ------------------------------------------------------------------------------
\subsection{Métriques et monitoring}
% ------------------------------------------------------------------------------

Le serveur expose un endpoint \texttt{/metrics} compatible Prometheus.
L'instrumentation est automatique via le package
\texttt{prometheus\_fastapi\_instrumentator}, complète par des métriques
personnalisées définies dans \texttt{metrics.py} :

\begin{itemize}
  \item \texttt{edr\_events\_ingested\_total} -- compteur d'événements
    par type ;
  \item \texttt{edr\_agents\_registered\_total} -- total des heartbeats ;
  \item \texttt{edr\_active\_agents} -- nombre d'agents actifs (heartbeat
    dans les 5 dernières minutes) ;
  \item \texttt{edr\_vt\_enrichments\_total} -- enrichissements VT
    (succès/échoués) ;
  \item \texttt{edr\_auth\_failures\_total} -- tentatives
    d'authentification échouées ;
  \item \texttt{edr\_fim\_events\_total} -- événements FIM par action ;
  \item \texttt{edr\_db\_health} -- état de la base de données
    (1=sain, 0=degrade).
\end{itemize}

Un collecteur personnalisé (\texttt{EDRDatabaseCollector}) interroge la
base de données à chaque \emph{scrape} Prometheus pour exposer des
jauges comme le nombre total d'événements stockés, la répartition par
type et sévérité, et les statistiques du cache VT.

% ------------------------------------------------------------------------------
\subsection{Dockerisation}
% ------------------------------------------------------------------------------

Le serveur est containerise via un Dockerfile multi-etapes. L'image est
publiée sur \texttt{ghcr.io/lockbits-esgi/edr:latest} avec support
multi-architecture (\texttt{linux/amd64} et \texttt{linux/arm64}).

Le conteneur expose le port 8000 et intègre un healthcheck HTTP sur
\texttt{/health}. La configuration s'effectue entièrement via des
variables d'environnement, ce qui permet un déploiement sans fichier de
configuration.

Un exemple d'exécution autonome :

\begin{lstlisting}[language=bash, caption={Execution du serveur EDR avec Docker}, label={lst:edr-docker}]
docker run -d \
  --name lockbits_edr \
  -p 8000:8000 \
  -e DATABASE_URL=sqlite:///./miniedr.db \
  -e AUTH_TOKEN=mon_token_secret \
  -e VT_ENABLED=false \
  -v edr_data:/app/data \
  ghcr.io/lockbits-esgi/edr:latest
\end{lstlisting}


% ==============================================================================
\section{EDR Agent}
% ==============================================================================

L'agent EDR est le composant installé sur les postes clients. Il collecte
les informations de sécurité et les transmet au serveur central.

% ------------------------------------------------------------------------------
\subsection{Architecture de l'agent}
% ------------------------------------------------------------------------------

L'agent est écrit en \textbf{Python 3.12} et compile en binaire autonome
avec \textbf{PyInstaller}. Cette approche évite toute dépendance
d'exécution sur la machine cible.

L'architecture est composee de quatre modules :

\begin{itemize}
  \item \texttt{agent/main.py} -- point d'entrée et logique métier
    (agent principal, modes scan et monitor) ;
  \item \texttt{agent/heartbeat.py} -- thread d'arrière-plan pour les
    heartbeats périodiques ;
  \item \texttt{agent/sender.py} -- client HTTP avec file d'attente
    locale pour la transmission fiable des événements ;
  \item \texttt{fim.py} -- surveillance des modifications de
    fichiers (File Integrity Monitoring) via watchdog ;
  \item \texttt{collector.py} -- collecte d'instantanés système
    (processus, connexions réseau, utilisateurs) ;
  \item \texttt{hasher.py} -- calcul d'empreintes SHA256 pour la
    détection de fichiers suspects ;
  \item \texttt{reporter.py} -- génération de rapports JSON et HTML
    pour l'export des données.
\end{itemize}

La configuration s'effectue via un fichier \texttt{config.yaml} ou par
variables d'environnement (\texttt{MINIEDR\_SERVER\_URL},
\texttt{MINIEDR\_AUTH\_TOKEN}, \texttt{MINIEDR\_AGENT\_COMPANY}).

\begin{lstlisting}[language=Python, caption={Structure de la classe MiniEDRAgent}, label={lst:agent-class}]
class MiniEDRAgent:
    """Agent that collects events and sends to server."""

    def __init__(self, config_path: Path, output_dir: Path,
                 user: str | None = None):
        self.config = _load_agent_config(config_path)
        self.agent_id = load_agent_id(...)
        self.company = self._load_company()
        self.sender = EventSender(
            server_url=server_url,
            auth_token=auth_token,
            timeout=server_config.get("timeout", 10),
            queue_path=queue_path,
        )
        self.heartbeat_worker: Optional[HeartbeatWorker] = None

    def create_event(self, event_type, severity, payload, tags=None):
        """Create MiniEDREvent with standard fields."""
        return MiniEDREvent(
            event_id=generate_uuid(),
            agent_id=self.agent_id,
            hostname=socket.gethostname(),
            platform=platform.system(),
            event_type=event_type,
            severity=severity,
            timestamp=utc_now_iso(),
            source="agent",
            payload=payload,
            tags=merge_company_tag(tags, self.company),
        )
\end{lstlisting}

% ------------------------------------------------------------------------------
\subsection{Fonctionnalités}
% ------------------------------------------------------------------------------

\paragraph{Mode scan.} Effectue un instantané ponctuel du système :
processus en cours (PID, nom, CPU, RAM), connexions réseau actives,
utilisateurs connectés, et informations système (OS, noyau, uptime).
Les données sont envoyées au serveur sous forme d'événement de type
\texttt{scan}.

\paragraph{Mode monitor.} Surveillance continue avec heartbeat périodique
(toutes les 5 minutes par défaut) et analyse FIM en temps réel.
L'agent surveille les répertoires configurés et génère des événements
\texttt{fim} lors de créations, modifications, suppressions ou
déplacements de fichiers. Un instantané système est également envoyé à
intervalles réguliers (toutes les 10 minutes par défaut).

\paragraph{FIM (File Integrity Monitoring).} La détection des
modifications de fichiers repose sur la bibliotheque \texttt{watchdog}.
Lorsqu'un fichier est créé ou modifié dans un répertoire surveillé,
l'agent calcule son empreinte SHA256 et l'inclut dans l'événement. Les
fichiers executables (.exe, .dll, .bat, .ps1, .sh) et les chemins
sensibles (/tmp/, /var/tmp/, startup) sont marques comme suspects.

\paragraph{Transmission fiable.} Le module \texttt{sender.py} implémente
un mécanisme de résilience : en cas d'indisponibilité du serveur, les
événements sont mis en file d'attente locale au format JSONL et
retransmis automatiquement lors du prochain \emph{flush}.

% ------------------------------------------------------------------------------
\subsection{Compilation PyInstaller}
% ------------------------------------------------------------------------------

La compilation en binaire autonome est réalisée par des scripts spécifiques
à chaque plateforme :

\begin{itemize}
  \item \texttt{packaging/build\_agent\_linux.sh} -- compilation pour
    Linux (amd64 et arm64) ;
  \item \texttt{packaging/build\_agent\_macos.sh} -- compilation pour
    macOS ;
  \item \texttt{packaging/build\_agent\_windows.bat} -- compilation pour
    Windows.
\end{itemize}

Les binaires sont distribues via les GitHub Releases du dépôt principal.
Chaque push sur \texttt{main} déclenche la CI qui compile les binaires
et met à jour les releases.

% ------------------------------------------------------------------------------
\subsection{Déploiement}
% ------------------------------------------------------------------------------

Le déploiement de l'agent s'effectue via le script
\texttt{install-agent.sh} qui supporte deux modes :

\begin{itemize}
  \item \textbf{Mode binaire} (recommandé) -- télécharge le binaire,
    l'installe dans \texttt{/opt/lockbits-agent/} et configure le service
    systemd ;
  \item \textbf{Mode Docker} -- tire l'image
    \texttt{ghcr.io/lockbits-esgi/edr-agent:latest} et l'exécute.
\end{itemize}

La commande d'installation typique est fournie directement sur le
dashboard client :

\begin{lstlisting}[language=bash, caption={Installation de l'agent EDR}, label={lst:agent-install}]
curl -fsSL https://raw.githubusercontent.com/Lockbits-ESGI/main/main/install-agent.sh | \
  SERVER_URL=https://edr.lockbits.pro AUTH_TOKEN=mon_token bash
\end{lstlisting}

L'agent supporte les commandes \texttt{--update} (mise à jour automatique)
et \texttt{--uninstall} (desinstallation complète).


% ==============================================================================
\section{Site Client}
% ==============================================================================

Le site client LockBits est l'interface web qui permet aux clients de
visualiser leur tableau de bord, gèrer leurs tickets de support, et
installer l'agent EDR.

% ------------------------------------------------------------------------------
\subsection{Architecture du site}
% ------------------------------------------------------------------------------

Le site est développé en \textbf{PHP 8+} avec \textbf{Apache} comme
serveur web. Il suit une architecture modulaire dans le répertoire
\texttt{site\_lockbits/client/} :

\begin{itemize}
  \item \texttt{auth.php} -- gestion des sessions et authentification ;
  \item \texttt{login.php}, \texttt{register.php} -- pages
    d'authentification ;
  \item \texttt{dashboard.php} -- tableau de bord client avec métriques
    EDR et commande d'installation de l'agent ;
  \item \texttt{tickets.php}, \texttt{ticket.php},
    \texttt{create\_ticket.php} -- gestion des tickets de support ;
  \item \texttt{chat\_client.php} -- interface du chatbot ;
  \item \texttt{config.php} -- configuration centralisée ;
  \item \texttt{db.php} -- connexion à la base de données via PDO.
\end{itemize}

% ------------------------------------------------------------------------------
\subsection{Modules fonctionnels}
% ------------------------------------------------------------------------------

\paragraph{Authentification.} Le module \texttt{auth.php} gère les
sessions PHP sécurisées avec régénération d'ID, cookie
\texttt{httponly} et \texttt{samesite=Lax}. Les mots de passe sont
hachés avec les algorithmes natifs de PHP.

\begin{lstlisting}[language=PHP, caption={Connexion à la base de données avec PDO}, label={lst:db-connection}]
function db(): PDO
{
    static $pdo = null;

    if ($pdo instanceof PDO) {
        return $pdo;
    }

    $dsn = sprintf('mysql:host=%s;port=%s;dbname=%s;charset=utf8mb4',
                   DB_HOST, DB_PORT, DB_NAME);

    $pdo = new PDO($dsn, DB_USER, DB_PASS, [
        PDO::ATTR_ERRMODE => PDO::ERRMODE_EXCEPTION,
        PDO::ATTR_DEFAULT_FETCH_MODE => PDO::FETCH_ASSOC,
        PDO::ATTR_EMULATE_PREPARES => false,
    ]);

    return $pdo;
}
\end{lstlisting}

\paragraph{Dashboard.} Le tableau de bord affiche des cartes de
statistiques (statut du service, uptime 30 jours, tickets ouverts, score
de sécurité) ainsi que la section de déploiement de l'agent EDR avec la
commande d'installation pre-remplie. L'interface est construite avec
Tailwind CSS et du JavaScript vanilla.

\paragraph{Tickets.} Les clients peuvent créer, visualiser et suivre
leurs tickets de support. Les tickets sont synchronisés avec GLPI pour
assurer la cohérence entre le portail client et l'outil interne.

% ------------------------------------------------------------------------------
\subsection{Base de données MySQL}
% ------------------------------------------------------------------------------

La base de données MySQL du site client contient les tables suivantes,
definies dans \texttt{client/database.sql} :

\begin{itemize}
  \item \texttt{users} -- comptes clients avec référence vers GLPI
    (\texttt{glpi\_user\_id}) ;
  \item \texttt{tickets} -- tickets de support relies aux utilisateurs
    et a GLPI (\texttt{glpi\_ticket\_id}) ;
  \item \texttt{ticket\_messages} -- messages associes aux tickets ;
  \item \texttt{organizations} -- organisations clients ;
  \item \texttt{edr\_endpoints} -- terminaux enregistrés par organisation ;
  \item \texttt{edr\_reports} -- rapports EDR reçus ;
  \item \texttt{edr\_report\_processes} -- processus collectés ;
  \item \texttt{edr\_report\_system} -- informations système.
\end{itemize}

% ------------------------------------------------------------------------------
\subsection{Integration GLPI SSO}
% ------------------------------------------------------------------------------

L'intégration GLPI (\texttt{glpi\_api.php}) permet l'authentification
unique (SSO) et la synchronisation des tickets. Elle supporte deux modes
de communication avec GLPI :

\begin{itemize}
  \item \textbf{API REST historique} (\texttt{apirest.php}) --
    authentification par \texttt{App-Token} et \texttt{Session-Token} ;
  \item \textbf{API REST v2} (\texttt{api.php}) -- authentification
    OAuth2 avec JWT.
\end{itemize}

Les fonctionnalités incluent la création d'utilisateurs GLPI lors de
l'inscription, la création de tickets depuis le portail, la
synchronisation des messages (followups), et la recupération de la liste
des tickets d'un utilisateur via l'API v2.

% ------------------------------------------------------------------------------
\subsection{Dockerisation}
% ------------------------------------------------------------------------------

L'image Docker du site client (\texttt{site\_lockbits/Dockerfile}) est
basee sur \texttt{php:8.2-apache}. Elle intégré :

\begin{itemize}
  \item Les extensions PHP nécessaires (mysqli, pdo\_mysql, curl, gd,
    zip, opcache) ;
  \item Le module Apache \texttt{rewrite} pour les URLs propres ;
  \item Le rate limiting Apache (200 req/min, burst de 50) ;
  \item Un script d'entrée (\texttt{docker-entrypoint.sh}) qui
    initialise le schéma de base de données au premier démarrage ;
  \item Un healthcheck HTTP sur la page d'accueil.
\end{itemize}

L'image est publiée sur \texttt{ghcr.io/lockbits-esgi/site\_lockbits:latest}.
Le déploiement en production nécessite un conteneur MySQL séparé (image
\texttt{mysql:8.0}), comme défini dans le fichier
\texttt{docker-compose.prod.yml}.


% ==============================================================================
\section{Chatbot}
% ==============================================================================

Le chatbot LockBits est un assistant virtuel alimente par l'API Claude
d'Anthropic. Il offre un support client interactif directement depuis
le portail client.

% ------------------------------------------------------------------------------
\subsection{Architecture}
% ------------------------------------------------------------------------------

Le chatbot est implemente dans un fichier PHP unique
(\texttt{site\_lockbits/chat.php}) qui sert de proxy entre le front-end
JavaScript et l'API Anthropic.

L'endpoint attend une requête POST avec un corps JSON contenant le
message de l'utilisateur et l'historique de la conversation. Il construit
un payload et l'envoie à l'API Claude :

\begin{lstlisting}[language=PHP, caption={Envoi d'une requête à l'API Claude}, label={lst:chat-payload}]
$payload = [
    'model'      => 'claude-sonnet-4-6',
    'max_tokens' => 500,
    'system'     => $systemPrompt,
    'messages'   => $messages,
];

$ch = curl_init('https://api.anthropic.com/v1/messages');
curl_setopt_array($ch, [
    CURLOPT_RETURNTRANSFER => true,
    CURLOPT_POST           => true,
    CURLOPT_POSTFIELDS     => json_encode($payload),
    CURLOPT_HTTPHEADER     => [
        'Content-Type: application/json',
        'x-api-key: ' . $apiKey,
        'anthropic-version: 2023-06-01',
    ],
]);
\end{lstlisting}

% ------------------------------------------------------------------------------
\subsection{Fonctionnalités}
% ------------------------------------------------------------------------------

L'assistant virtuel est capable de :

\begin{itemize}
  \item     Répondre aux questions sur les services LockBits (hébergement,
    cybersécurité, infrastructure) ;
  \item Fournir un resume du tableau de bord client (tentatives DDoS
    bloquees, vulnerabilites, incidents) ;
  \item Conserver l'historique de la conversation entre chaque message ;
\item S'adapter à la langue du client (français par défaut, anglais
si nécessaire).
\end{itemize}

Le \texttt{system prompt} contient les informations essentielles sur
l'entreprise : SLA de 99,99\%, monitoring SOC 24/7/365, temps de réponse
aux incidents inférieur à 5 minutes, certification ISO 27001 Ready.

% ------------------------------------------------------------------------------
\subsection{Sécurité}
% ------------------------------------------------------------------------------

Plusieurs mesures de sécurité encadrent le chatbot :

\begin{itemize}
  \item \textbf{Rate limiting} -- compteur quotidien base sur un fichier
    \texttt{usage\_count.txt} qui limite à 100 requêtes par jour toutes
    sources confondues. Le compteur est reinitialise chaque jour ;
  \item \textbf{Validation des entrées} -- la méthode POST est obligatoire,
    et le message est vérifié avant traitement ;
  \item \textbf{Clé API} -- la clé d'API Claude est chargée depuis la
    variable d'environnement \texttt{CLAUDE\_API\_KEY}, jamais exposée
    dans le code source ;
  \item     \textbf{Périmètre restreint} -- le prompt système interdit au
    chatbot de répondre à des questions hors du domaine de LockBits et
    de la cybersécurité.
\end{itemize}

L'interface utilisateur du chatbot est integree directement dans le
dashboard client via du JavaScript vanilla. Elle affiche un badge de
notification en cas de nouveau message non lu et propose une experience
de messagerie instantanée avec indicateur de saisie.
